\section{Experimental Design}\label{sec:experiments}
\paragraph{Evaluated models.} We compare \Micro{}, \Macro{}, and \Hybrid{} with four implemented baselines: WFlib DF, WFlib VarCNN, a 104-dimensional CUMUL-style RBF-SVM, and WFlib-based Tik-Tok. Contrastive NetCLR is the fifth planned SOTA baseline; its outcomes are not available at the time of this manuscript version.

\paragraph{Baseline fitting.} DF and VarCNN were fitted to Historical DEV, with DEV\_LATE used for selection and phase-C refit performed after fixing the selected training length. CUMUL-style uses 104 traffic statistics and an RBF SVM with $C=1$ and scale-dependent $\gamma$. Its packet-size extraction may include protocol-header overhead and should not be called an exact reproduction of the original CUMUL method. Tik-Tok consumes directional timestamps of length 5,000 and uses the fixed WFlib architecture, Adamax at learning rate $2\times10^{-3}$, and an epoch selected on DEV\_LATE by macro-F1. Its phase-B best epoch was 23, followed by a 23-epoch phase-C refit over all Historical DEV observations.

\paragraph{Metrics.} We report top-1 accuracy, macro-averaged F1, top-5 accuracy, and mean reciprocal rank (MRR) computed from the frozen $K=65$ score vector. Macro-F1 reflects site-level imbalance; MRR reflects the position of the correct site in the complete ranking. Values are computed independently within each cohort.

\paragraph{Scope and reproducibility.} The same aligned capture identifiers and labels are required for every model in a cohort. The published values are single-run results, not claims of cross-seed stability. Hardware, environment versions, training duration, and architecture provenance should be documented for each method. NetCLR will be reported as WFlib-adapted because its practical contrastive recipe differs from the upstream batch size and dataset regimen.
